\begin{center}\itshape
Take a polynomial at random. It is irreducible, its roots are as unrelated as roots can be,\\
and modulo a prime it breaks like a shuffled deck of cards.
\end{center}

In 1916 Hermann Weyl published in the \emph{Mathematische Annalen} a long paper on the distribution of numbers modulo one, whose criterion (Volume~5, Theorem~8.13) turned ``random-looking'' into a theorem about exponential sums. Thirty years later Paul Erd\H{o}s and Paul Turán made it quantitative: the discrepancy of a finite sequence, its largest deviation from uniformity on an interval, is controlled by finitely many of Weyl's sums. This chapter starts with their inequality and with its cousin on a compact group, which makes sense of the Sato--Tate law of Volume~13: the angles of the Frobenius of elliptic curves follow the density $\frac2\pi\sin^2\theta$, the law of the eigenvalue angle of a random rotation of $SU(2)$.

The second half of the chapter is about polynomials, and the randomness is that of permutations. In 1936 Bartel van der Waerden, then in Leipzig and the author of the most influential algebra textbook of the century, proved that a random monic integer polynomial of degree $n$ has, with probability tending to $1$, the largest possible Galois group, the symmetric group $S_n$. His title said it all: \emph{Die Seltenheit der reduziblen Gleichungen und der Gleichungen mit Affekt}, the rarity of reducible equations and of equations with ``affect'' (a smaller Galois group). He conjectured that the exceptions are dominated by the reducible polynomials, and the conjecture stood for 85 years, until Manjul Bhargava proved it in a preprint of 2021.

Two theorems drive the proof we give. First, modulo a fixed prime, a random polynomial factors like a random permutation decomposes into cycles: the probability of each factorization pattern is Cauchy's probability for the corresponding cycle type. Second, a theorem of Camille Jordan from 1872: a proper subgroup of a finite group misses an entire conjugacy class. Dedekind's theorem (Volume~13) links the two, and a variance argument in the spirit of Turán's proof of the Hardy--Ramanujan theorem does the counting.

\begin{reminder}[from Volumes~3, 5 and~13]
\emph{Irreducible polynomials} (Volume~3, Theorem~9.10). The number $N_p(d)$ of monic irreducible polynomials of degree $d$ over $\mathbb F_p$ satisfies $\sum_{e\mid d}eN_p(e)=p^d$, hence $N_p(d)=\frac1d\sum_{e\mid d}\mu(e)p^{d/e}$.
\emph{Weyl's criterion} (Volume~5, Theorem~8.13). A sequence $(x_n)$ in $\R/\Z$ is equidistributed if and only if $\frac1N\sum_{n\le N}e(kx_n)\to0$ for every integer $k\ne0$, where $e(x)=e^{2\pi ix}$.
\emph{Dedekind's theorem} (Volume~13, Theorem~4.20). Let $f\in\Z[X]$ be monic and irreducible of degree $n$, with Galois group $G\subseteq S_n$ acting on its roots, and $p\nmid\operatorname{disc}(f)$. If $\bar f$ factors in $\mathbb F_p[X]$ into irreducibles of degrees $d_1,\dots,d_r$, then $G$ contains a permutation of cycle type $(d_1,\dots,d_r)$.
\emph{Chebotarev} (Volume~13, Theorem~11.23). With $f$ as above, the primes $p$ for which $\bar f$ has a given factorization type $\lambda$ have density $\frac{\#\{g\in G:g\text{ of type }\lambda\}}{|G|}$.
\emph{Hasse} (Volume~13, Corollary~16.14). For $p$ odd and $x^3+ax+b$ without repeated root mod $p$, the elliptic curve $y^2=x^3+ax+b$ has $|E(\mathbb F_p)|=p+1-a_p$ with $|a_p|\le2\sqrt p$.
\end{reminder}

\section{Discrepancy}

For points $x_1,\dots,x_N$ of $\R/\Z$, the \emph{discrepancy} is
\[
D_N=\sup_I\Bigl|\frac{\#\{n\le N:x_n\in I\}}N-|I|\Bigr|,
\]
the supremum over all arcs $I$ of the circle. A sequence is equidistributed exactly when $D_N\to0$. Write $S_k=\sum_{n\le N}e(kx_n)$.

\begin{theorem}[an Erd\H{o}s--Turán inequality]\label{thm:erdosturan}
For every integer $K\ge2$,
\[
D_N\le\frac6{\sqrt K}+\frac2\pi\sum_{k=1}^{K-1}\frac{|S_k|}{kN}.
\]
\end{theorem}
\begin{proof}
Let $F_K(t)=\sum_{|k|<K}(1-\frac{|k|}K)e(kt)=\frac1K\bigl(\frac{\sin\pi Kt}{\sin\pi t}\bigr)^2$ be Fejér's kernel: $F_K\ge0$ and $\int_0^1F_K=1$. Since $\sin\pi t\ge2t$ on $[0,\frac12]$, $F_K(t)\le\frac1{4Kt^2}$ for $0<|t|\le\frac12$, so with $\delta=K^{-1/2}$ the mass of $F_K$ outside $[-\delta,\delta]$ is $\varepsilon\le2\int_\delta^\infty\frac{\dd t}{4Kt^2}=\frac1{2\sqrt K}$. Let $E$ be the sum over $k$ in the statement; we may assume $E\le1$, since $D_N\le1$.

For an arc $J$ let $g_J=\mathbf 1_J*F_K$, a trigonometric polynomial with $\hat g_J(0)=|J|$ and $|\hat g_J(k)|\le|\hat{\mathbf 1}_J(k)|=\frac{|\sin\pi k|J||}{\pi|k|}\le\frac1{\pi|k|}$ for $0<|k|<K$. Then
\[
\Bigl|\frac1N\sum_ng_J(x_n)-|J|\Bigr|=\Bigl|\sum_{0<|k|<K}\hat g_J(k)\frac{S_k}N\Bigr|\le E,
\]
because $|S_{-k}|=|S_k|$.

\emph{Upper bound.} Let $I$ be an arc of length $\ell$, and $J$ the arc of length $\ell+2\delta$ with the same centre (if $\ell+2\delta\ge1$ the upper bound below is trivial). For $x\in I$, $g_J(x)=\int\mathbf 1_J(x-t)F_K(t)\dd t\ge\int_{|t|\le\delta}F_K=1-\varepsilon$. Hence, as $\varepsilon\le\frac12$,
\[
\frac{\#\{n:x_n\in I\}}N\le\frac{\ell+2\delta+E}{1-\varepsilon}\le(\ell+2\delta+E)(1+2\varepsilon)\le\ell+2\delta+E+8\varepsilon\le\ell+\frac6{\sqrt K}+E.
\]

\emph{Lower bound.} If $\ell\le2\delta$ there is nothing to prove. Otherwise let $J$ be the arc of length $\ell-2\delta$ with the same centre. For $x\notin I$, every point of $J$ is at distance more than $\delta$ from $x$, so $g_J(x)\le\varepsilon$; and $g_J\le1$ everywhere. So $g_J\le\mathbf 1_I+\varepsilon$, and $\frac{\#\{n:x_n\in I\}}N\ge\frac1N\sum_ng_J(x_n)-\varepsilon\ge\ell-2\delta-E-\varepsilon\ge\ell-\frac3{\sqrt K}-E$.
\end{proof}

Erd\H{o}s and Turán proved the inequality with $\frac CK$ in place of $\frac6{\sqrt K}$; the sharp form needs trigonometric polynomials that majorize an interval with an error $O(\frac1K)$, which Selberg and Vaaler constructed later. For most purposes the shape is what matters: \emph{finitely many Weyl sums control the discrepancy}, and Weyl's criterion becomes a rate (Exercise~\ref{ex:golden}).

\section{The Sato--Tate measure}

The same principle works on other compact spaces, with characters in place of $e(kx)$. We need one case. The \emph{Sato--Tate measure} on $[0,\pi]$ is $\mu_{ST}=\frac2\pi\sin^2\theta\dd\theta$. It is the law of the angle $\theta$ of the eigenvalues $e^{\pm i\theta}$ of a random matrix of $SU(2)$, chosen with Haar measure, by Weyl's integration formula (Volume~9, Theorem~8.8), and its image under $\theta\mapsto x=\cos\theta$ is the semicircle law $\frac2\pi\sqrt{1-x^2}\dd x$. Let
\[
U_k(\theta)=\frac{\sin(k+1)\theta}{\sin\theta}=\sum_{j=0}^ke^{i(k-2j)\theta},
\]
which is a polynomial of degree $k$ in $\cos\theta$, with leading coefficient $2^k$ (the Chebyshev polynomial of the second kind); it is the trace of the $k$-th symmetric power of the matrix.

\begin{proposition}\label{prop:stcriterion}
(a) $\int U_jU_k\dd\mu_{ST}=\mathbf 1_{j=k}$ for $j,k\ge0$. (b) Probability measures $\nu_N$ on $[0,\pi]$ converge weakly to $\mu_{ST}$ if and only if $\int U_k\dd\nu_N\to0$ for every $k\ge1$.
\end{proposition}
\begin{proof}
(a) $\int_0^\pi\frac2\pi\sin(j+1)\theta\sin(k+1)\theta\dd\theta=\mathbf 1_{j=k}$. (b) If $\nu_N\to\mu_{ST}$ weakly, then $\int U_k\dd\nu_N\to\int U_k\dd\mu_{ST}=\int U_kU_0\dd\mu_{ST}=0$, as $U_k$ is continuous. Conversely, the $U_k$, $k\le m$, span the polynomials of degree $\le m$ in $\cos\theta$, since their degrees are exactly $0,1,\dots,m$. Let $\varphi$ be continuous on $[0,\pi]$ and $\eta>0$. As $\cos$ is a homeomorphism of $[0,\pi]$ onto $[-1,1]$, the Weierstrass theorem gives a polynomial $Q$ with $|\varphi(\theta)-Q(\cos\theta)|\le\eta$ for all $\theta$. Writing $Q(\cos\theta)=\sum_{k\le m}c_kU_k(\theta)$, $\int Q(\cos\theta)\dd\nu_N\to c_0=\int Q(\cos\theta)\dd\mu_{ST}$, so $\limsup|\int\varphi\dd\nu_N-\int\varphi\dd\mu_{ST}|\le2\eta$.
\end{proof}

For an elliptic curve $E$ over $\Q$ without complex multiplication, write $a_p=2\sqrt p\cos\theta_p$ with $\theta_p\in[0,\pi]$ (Hasse). The Sato--Tate conjecture, now a theorem of Laurent Clozel, Michael Harris, Nicholas Shepherd-Barron and Richard Taylor, completed with Thomas Barnet-Lamb and David Geraghty in 2011, says that the $\theta_p$ are equidistributed for $\mu_{ST}$ as $p$ runs over the primes (Volume~13 shows the histogram for $y^2=x^3-x+1$). By Proposition~\ref{prop:stcriterion} this is the statement $\frac1{\pi(x)}\sum_{p\le x}U_k(\theta_p)\to0$ for each $k\ge1$, a prime number theorem for the $L$-functions of the symmetric powers of $E$, and that is how it was proved.

The \emph{vertical} version, where the prime is fixed and the curve varies, is elementary. Let $\chi$ be the Legendre symbol modulo $p\ge5$ and, for $a,b\in\mathbb F_p$, $a_p(a,b)=-\sum_{x\in\mathbb F_p}\chi(x^3+ax+b)$; when $\Delta=4a^3+27b^2\ne0$ this is the $a_p$ of the curve $E_{a,b}:y^2=x^3+ax+b$, since each $x$ contributes $1+\chi(x^3+ax+b)$ affine points.

\begin{proposition}[Birch, 1968]\label{prop:birch}
For a prime $p\ge5$,
\[
\sum_{a,b\in\mathbb F_p}a_p(a,b)^2=p^2(p-1),\qquad\frac1{\#\{(a,b):\Delta\ne0\}}\sum_{\Delta\ne0}\Bigl(\frac{a_p(a,b)}{\sqrt p}\Bigr)^2=1+O\Bigl(\frac1p\Bigr).
\]
\end{proposition}
\begin{proof}
Expanding the square, $\sum_{a,b}a_p^2=\sum_{x,y}\sum_{a,b}\chi(x^3+ax+b)\chi(y^3+ay+b)$. If $x=y$, the inner sum counts the $(a,b)$ with $x^3+ax+b\ne0$: there are $p^2-p$, and the $p$ values of $x$ give $p(p^2-p)$. If $x\ne y$, the map $(a,b)\mapsto(u,v)=(x^3+ax+b,y^3+ay+b)$ is a bijection of $\mathbb F_p^2$ (its linear part has determinant $x-y\ne0$), so the inner sum is $(\sum_u\chi(u))^2=0$. This proves the identity.

There are $p$ pairs with $\Delta=0$, namely $(-3t^2,2t^3)$ for $t\in\mathbb F_p$ (if $\Delta=0$ and $a\ne0$, then $t=-\frac{3b}{2a}$ works, and $a=0$ forces $b=0$). For them $x^3+ax+b=(x-r)^2(x-s)$, and $a_p=-\sum_{x\ne r}\chi(x-s)=\chi(r-s)$, so $|a_p|\le1$. Hence $\sum_{\Delta\ne0}a_p^2=p^3-p^2+O(p)$, and dividing by $p(p^2-p)$ gives the second claim.
\end{proof}

So on average over curves, $(2\cos\theta)^2=\frac{a_p^2}p$ has mean $1$, which is $\int(2\cos\theta)^2\dd\mu_{ST}=\int(U_2+1)\dd\mu_{ST}$. Birch computed all the moments in the same way, with more effort: the $2m$-th moment tends to the Catalan number $C_m=\frac1{m+1}\binom{2m}m$, the $2m$-th moment of $2\cos\theta$ under $\mu_{ST}$ (Exercise~\ref{ex:catalan}). Figure~\ref{fig:birch} shows the vertical law for $p=1009$.

\begin{figure}[ht]
\centering\fig{ch12-birch}
\caption{All $1\,017\,072$ elliptic curves $y^2=x^3+ax+b$ over $\mathbb F_{1009}$: histogram of $\frac{a_p}{2\sqrt p}$ against the semicircle density $\frac2\pi\sqrt{1-x^2}$. The moments of $\frac{a_p}{\sqrt p}$ of orders $2$, $4$, $6$ are $1.0000$, $2.0000$, $5.0000$: the Catalan numbers $1$, $2$, $5$.}
\label{fig:birch}
\end{figure}

\section{Random polynomials modulo $p$}

A monic polynomial $f$ of degree $n$ over $\mathbb F_p$ without repeated factor has a \emph{factorization type}, the partition $\lambda$ of $n$ given by the degrees of its irreducible factors. Write $\lambda=1^{m_1}2^{m_2}\cdots n^{m_n}$, with $m_d$ factors of degree $d$ and $\sum_ddm_d=n$. A permutation of $\{1,\dots,n\}$ has a \emph{cycle type}, also a partition of $n$, and by Cauchy's formula the number of permutations of type $\lambda$ is $\frac{n!}{\prod_dd^{m_d}m_d!}$ (choose the order of the $n$ symbols, then forget the $d$ rotations of each $d$-cycle and the $m_d!$ orders of the cycles of length $d$). Let $P(\lambda)=\prod_d\frac1{d^{m_d}m_d!}$ be the probability that a uniform random permutation has type $\lambda$.

\begin{theorem}\label{thm:factortypes}
Let $N_\lambda(p)$ be the number of monic squarefree polynomials of degree $n$ over $\mathbb F_p$ with factorization type $\lambda$. Then
\[
N_\lambda(p)=\prod_d\binom{N_p(d)}{m_d}=p^nP(\lambda)\Bigl(1+O_n\Bigl(\frac1p\Bigr)\Bigr).
\]
Moreover, at most $2p^{n-1}$ monic polynomials of degree $n$ are not squarefree.
\end{theorem}
\begin{proof}
A squarefree polynomial of type $\lambda$ is the product of a set of $m_d$ distinct monic irreducibles of degree $d$, for each $d$, and the factorization is unique; this gives the product of binomial coefficients. By Gauss's formula, $|dN_p(d)-p^d|\le\sum_{e\le d/2}p^e\le2p^{d/2}$, so $N_p(d)=\frac{p^d}d(1+O(p^{-d/2}))$ for $d\ge2$ and $N_p(1)=p$. For fixed $m$, $\binom Nm=\frac{N^m}{m!}(1+O(\frac{m^2}N))$. Hence $\binom{N_p(d)}{m_d}=\frac{p^{dm_d}}{d^{m_d}m_d!}(1+O_n(\frac1p))$, and the product over $d$ is $p^nP(\lambda)(1+O_n(\frac1p))$. A polynomial that is not squarefree is $g^2h$ with $g$ monic of degree $k\ge1$: at most $\sum_{k\ge1}p^kp^{n-2k}\le2p^{n-1}$ choices.
\end{proof}

So, for $p$ large, a random monic polynomial of degree $n$ over $\mathbb F_p$ is squarefree with probability $1-O(\frac1p)$, and its factorization type has the law of the cycle type of a random permutation of $S_n$. It is irreducible with probability about $\frac1n$ (an $n$-cycle), its number of roots in $\mathbb F_p$ is distributed like the number of fixed points of a random permutation, about Poisson of mean $1$ for large $n$, and the probability that it has no root at all is close to $\frac1e$ (Exercise~\ref{ex:noroot}).

\begin{think}
Why should a polynomial know about permutations? Over $\mathbb F_{p^{n!}}$, a squarefree $f$ of degree $n$ splits into distinct roots, and the Frobenius $x\mapsto x^p$ permutes them. What is the cycle type of this permutation?
\end{think}

\section{Galois groups of random polynomials}

\begin{lemma}[Jordan, 1872]\label{lem:jordan}
If $H$ is a proper subgroup of a finite group $G$, then $\bigcup_{g\in G}gHg^{-1}\ne G$. Equivalently, some conjugacy class of $G$ does not meet $H$.
\end{lemma}
\begin{proof}
The conjugate $gHg^{-1}$ depends only on the coset $gH$, so there are at most $[G:H]$ distinct conjugates. Each contains the identity, so their union has at most $[G:H](|H|-1)+1=|G|-[G:H]+1<|G|$ elements, as $[G:H]\ge2$.
\end{proof}

For an infinite group the lemma fails: every matrix of $GL_2(\mathbb C)$ is conjugate to an upper triangular one (Exercise~\ref{ex:jordaninfinite}).

For $H\ge1$ let $\mathcal B_n(H)$ be the set of monic $f=X^n+a_{n-1}X^{n-1}+\dots+a_0\in\Z[X]$ with $|a_i|\le H$; it has $V=(2H+1)^n$ elements. Call $f$ \emph{exceptional} if it is reducible over $\Q$, or irreducible with Galois group $G_f\ne S_n$.

\begin{lemma}\label{lem:missingtype}
If $f\in\Z[X]$ is monic of degree $n$ and exceptional, there is a partition $\lambda$ of $n$ such that $\bar f$ does not have factorization type $\lambda$ for any prime $p$ for which $\bar f\in\mathbb F_p[X]$ is squarefree.
\end{lemma}
\begin{proof}
If $f$ is reducible over $\Q$, then by Gauss's lemma $f=gh$ with $g,h\in\Z[X]$ monic of positive degrees, so $\bar f=\bar g\bar h$ is never irreducible: take $\lambda=(n)$. If $f$ is irreducible and $G_f\ne S_n$, Lemma~\ref{lem:jordan} gives a conjugacy class of $S_n$, that is a cycle type $\lambda$, containing no element of $G_f$. If $\bar f$ is squarefree, its discriminant, which is $\operatorname{disc}(f)\bmod p$, is non-zero, so $p\nmid\operatorname{disc}(f)$, and by Dedekind's theorem the type of $\bar f$ is the cycle type of an element of $G_f$, hence is not $\lambda$.
\end{proof}

\begin{theorem}[van der Waerden, 1936]\label{thm:vdw}
For each $n\ge2$ there is $C_n$ such that the number of exceptional $f\in\mathcal B_n(H)$ is at most $C_n(2H+1)^nH^{-1/3}\log H$ for $H\ge2$. In particular, almost all monic integer polynomials of degree $n$ are irreducible with Galois group $S_n$.
\end{theorem}
\begin{proof}
Fix a partition $\lambda$ of $n$; since there are finitely many, it suffices to bound the number of $f\in\mathcal B_n(H)$ for which no prime $p\le Q$ makes $\bar f$ squarefree of type $\lambda$, where $Q=H^{1/3}$. Let $\Omega_p\subseteq\mathbb F_p[X]$ be the set of monic squarefree polynomials of degree $n$ and type $\lambda$, $\rho_p=\frac{|\Omega_p|}{p^n}$, and $p_0=p_0(n)$ such that $\rho_p\ge\frac{P(\lambda)}2\ge\frac1{2\cdot n!}$ for $p>p_0$ (Theorem~\ref{thm:factortypes}). Let $\mathcal P$ be the set of primes in $]p_0,Q]$, $A=\sum_{p\in\mathcal P}\rho_p$, and $\omega(f)=\#\{p\in\mathcal P:\bar f\in\Omega_p\}$.

\emph{Counting in residue classes.} For a squarefree $m$ and a vector $r$ modulo $m$, the number of $f\in\mathcal B_n(H)$ with coefficients $\equiv r\pmod m$ is $\prod_{i<n}(\frac{2H+1}m+\theta_i)$ with $|\theta_i|\le1$, which is $\frac V{m^n}(1+\theta\eta_m)$ with $|\theta|\le1$ and $\eta_m=\frac{2nm}{2H+1}$, as long as $\eta_m\le1$. By the Chinese remainder theorem, the number of $f$ with $\bar f\in\Omega_p$ is $V\rho_p(1+\theta\eta_p)$, and for $p\ne q$ the number with $\bar f\in\Omega_p$ and $\bar f\in\Omega_q$ is $V\rho_p\rho_q(1+\theta\eta_{pq})$. All moduli used are at most $Q^2$; let $\eta=\frac{2nQ^2}{2H+1}\le\frac{nQ^2}H$.

\emph{The variance.} Summing over $f\in\mathcal B_n(H)$, $\sum\omega\ge VA(1-\eta)$, $\sum\omega\le VA(1+\eta)$ and $\sum\omega^2=\sum\omega+\sum_{p\ne q}V\rho_p\rho_q(1+\theta\eta_{pq})\le VA(1+\eta)+VA^2(1+\eta)$. Hence
\[
\sum_f(\omega(f)-A)^2=\sum\omega^2-2A\sum\omega+VA^2\le VA(1+\eta)+3\eta VA^2.
\]
An $f$ with $\omega(f)=0$ contributes $A^2$ to the left side, so the number of such $f$ is at most $\frac{V(1+\eta)}A+3\eta V$.

\emph{Conclusion.} With $Q=H^{1/3}$, $\eta\le nH^{-1/3}$, and by the prime number theorem $A\ge\frac{\pi(Q)-\pi(p_0)}{2\cdot n!}\ge\frac{Q}{3\cdot n!\log Q}$ for $H$ large. So the number of $f$ with $\omega(f)=0$ is at most $V(9\cdot n!\,H^{-1/3}\log H+3nH^{-1/3})$. By Lemma~\ref{lem:missingtype}, every exceptional $f$ has $\omega(f)=0$ for some $\lambda$, which proves the theorem for $H$ large, and for small $H$ by adjusting $C_n$.
\end{proof}

\begin{remark}
These are quoted for culture; see the Sources. (a) Patrick Gallagher (1973) replaced the variance by the large sieve and obtained the bound $O(H^{n-1/2}\log H)$. (b) The reducible polynomials alone number $\asymp H^{n-1}$ (for instance those with $a_0=0$), and van der Waerden conjectured that the exceptional ones are $O(H^{n-1})$. Bhargava proved this in 2021, with a combination of the geometry of numbers, the counting of polynomials with large discriminant divisors, and a careful study of each possible Galois group. (c) Hilbert's irreducibility theorem (1892) says that specializing the coefficients of the generic polynomial, whose Galois group is $S_n$, keeps the group $S_n$ for infinitely many integer values; van der Waerden's theorem says ``for almost all''.
\end{remark}

For a fixed $f$, Chebotarev's theorem says more: the types of $\bar f$, as $p$ varies, follow the law of the cycle type of a \emph{random element of $G_f$}. Figure~\ref{fig:chebotarev} compares a polynomial with group $S_5$ and one with a group of order $20$. Lemma~\ref{lem:missingtype} is visible there: the second polynomial never has the types $1^32$, $1^23$ or $23$.

\begin{figure}[ht]
\centering\fig{ch12-chebotarev}
\caption{Factorization types of $x^5-x-1$ (charcoal bars, Galois group $S_5$) and of $x^5-2$ (rust bars, Galois group the Frobenius group of order $20$) modulo the primes $7\le p<10^5$, against the proportions of the cycle types in each group (dots). For $x^5-2$ only four types occur, with frequencies close to $\frac1{20}$, $\frac5{20}$, $\frac{10}{20}$, $\frac4{20}$.}
\label{fig:chebotarev}
\end{figure}

\begin{beauty}
Three kinds of randomness meet in this chapter and turn out to be one. A random polynomial over $\mathbb F_p$ factors like a random permutation. A random integer polynomial has the full symmetric group, so, by Chebotarev, its reductions modulo the primes again factor like random permutations. And an elliptic curve, whose ``Galois group'' is a compact group rather than a finite one, has Frobenius angles that look like those of a random rotation. In each case the arithmetic object hides a group, and the statistics of the object are the statistics of a random element of the group.
\end{beauty}
